\begin{exercise}{59}
Consider the set $S=\{\partial\Delta^n\to\Delta^n\}_{n\ge0}$ given by the boundary inclusions. Show
that $\chi(S)$ is given by all monomorphisms of simplicial sets.

\begin{proof}
The nondegenerate simplices of $\partial\Delta^n$ correspond
to the nonempty proper subsets of $[n]$, so there are finitely
many. Lemma~1.3.15 therefore identifies $\chi(S)$ with the
saturated closure $\overline S$.
Monomorphisms form a saturated class containing $S$
(the example after Remark~1.3.7), so
$\overline S$ consists of monomorphisms.

Conversely, identify a monomorphism with an inclusion $A\subseteq B$.
Put $A_{-1}=A$ and $A_n=A\cup\mathrm{sk}_n(B)$ for $n\ge0$,
and let $J_n$ be the set of nondegenerate $n$-simplices of $B$
that do not belong to $A$.
As in Exercise~\ref{ex:15}, there is a pushout
\[
 \begin{tikzcd}[column sep=large,row sep=large]
  \displaystyle\coprod_{\sigma\in J_n}\partial\Delta^n
    \arrow[r] \arrow[d] & A_{n-1}\arrow[d] \\
  \displaystyle\coprod_{\sigma\in J_n}\Delta^n
    \arrow[r] & A_n .
 \end{tikzcd}
\]
Indeed, by Exercise~\ref{ex:5}, membership in a simplicial
subset is equivalent to membership of the nondegenerate term:
one recovers that term by applying a section of the degeneracy.
Thus, in each degree, the new simplices are precisely the
pullbacks of elements of $J_n$ along surjections, each occurring
once. Nonsurjective simplex operators give boundary simplices,
which already lie in $A_{n-1}$.
This is exactly the degreewise pushout above.
For $n=0$, it attaches the missing vertices along
$\partial\Delta^0=\varnothing$.

Every $A_{n-1}\to A_n$ is a pushout of a coproduct of elements
of $S$, hence lies in $\overline S$.
By Exercise~\ref{ex:15}, the union of the $A_n$ is $B$.
Closure under countable compositions now gives
$A\hookrightarrow B\in\overline S$.
Finally, an arbitrary monomorphism is isomorphic as an arrow
to this inclusion, hence is a retract of it and also belongs
to $\overline S$.
\end{proof}
\end{exercise}
